\section{Comparison with Binius64}
\label{sec:cmp}

We close by positioning the construction against \binius{}
\cite{binius64repo}, the SHA-256 prover of the Binius line
\cite{binius,fribinius} and the closest system in the literature. The two are
close cousins: \zinc{} imports \binius{}'s per-step modular-addition carry
identity (\Cref{sec:arith}) and its single-$\AND$ form of $\Ch$, and both work
in characteristic two precisely to avoid the prime-field witness blow-up that
bit-decomposing SHA-256 would incur. They differ in \emph{one} structural
choice --- \binius{} arithmetizes SHA-256 as an unrolled, non-deterministic
$64$-bit \emph{word circuit}, whereas \zinc{} uses a \emph{uniform} AIR over
the cell ring $\cellring$ --- and that choice, together with two independent
choices of commitment and reduction field, accounts cleanly for the four
numbers a user observes. This section states the three axes and the measured
outcome; the supporting anatomy (the shift handling, the AIR counterfactual,
deployment guidance, and the full measurement sweeps) is deferred to the
appendices.

\subsection{Three orthogonal design axes}
\label{sec:axes}

The comparison is unusually clean: the systems make independent choices on
three axes, and each axis decides exactly one observable outcome
(\Cref{tab:axes}).

\begin{table}[t]
\centering
\renewcommand{\arraystretch}{1.25}
\begin{tabularx}{\textwidth}{@{}lXXl@{}}
\toprule
\textbf{Axis} & \textbf{\binius{}} & \textbf{\zinc{}} & \textbf{Decides} \\
\midrule
Arithmetization &
  unrolled word \emph{circuit} (arbitrary wiring) &
  uniform \emph{AIR} over $\Ftwo[X]$ (structured wiring) &
  \textbf{verifier} \\
Commitment / proximity &
  Reed--Solomon $+$ recursive FRI/BaseFold &
  linear $\Ftwo$-RAA $+$ flat Brakedown &
  \textbf{proof size} \\
Reduction field &
  sumcheck tensors in $\mathrm{GF}(2^{128})$ &
  discharge kept in $\Ftwo/\mathrm{GF}(2^8)$ &
  \textbf{memory / large $N$} \\
\bottomrule
\end{tabularx}
\caption{The three design axes. Each system picks a different point on each,
and each axis decides one headline outcome. No single axis decides the prover:
it is dominated by content both systems share (the $\AND$ discharge and the
commitment encode), so the matched-rate prover is a net of offsetting per-phase
gaps that leaves \zinc{} a touch ahead and widening with $N$.}
\label{tab:axes}
\end{table}

\paragraph{Arithmetization decides the verifier.} A verifier is succinct
exactly when the constraint system's wiring --- the map from each constraint to
the witness positions it reads --- is succinctly evaluable at a random point.
\zinc{}'s uniform AIR has \emph{structured} wiring: one fixed transition
relation, evaluated as an $O(\log N)$ formula, with the verifier never reading
the per-row structure. \binius{}'s word circuit has \emph{arbitrary} wiring:
which committed word feeds which gate is data with no closed form, so the
verifier sums over all $N$ gates in the clear --- its $O(N)$ verifier, and, by
the same token, its $O(N)$ preprocessing against \zinc{}'s $O(1)$.
Crucially, it is the cross-entry \emph{wiring}, not the entry-wise shifts, that
is linear: the entry-wise shift is succinct in both systems
(\Cref{sec:shifts}), and whether \binius{} could recover a succinct verifier by
adopting an AIR is the counterfactual of \Cref{sec:ana-air}.

\paragraph{Commitment decides proof size.} Independently of the
arithmetization, \zinc{} buys a linear-time $\Ftwo$-RAA encoder and pays with a
\emph{flat} proof --- a Brakedown test opening $t\approx 987$ full columns, so
the proof grows as $t\cdot\sqrt{N}$ with a large floor --- while \binius{} buys
a \emph{small, recursive} FRI proof and pays with an NTT encoder. A WHIR-style
recursive proximity test would close this gap on the \zinc{} side
(\Cref{sec:cmp-outlook}).

\paragraph{Reduction field decides memory and scaling.} \binius{}'s BitAnd and
Shift reductions build $\mathrm{GF}(2^{128})$ tensors that, with a
rate-$\tfrac12$ codeword, spill cache super-linearly; \zinc{} keeps the witness
bit-packed and runs the Hadamard discharge in the base field
$\Ftwo/\mathrm{GF}(2^8)$, lifting to $\K$ only at the $\psi$ projections ---
roughly $16$--$32\times$ fewer bytes streamed and about half the footprint at
the largest matched instance.

\paragraph{No single axis decides the prover.} The prover is dominated by work
both systems must do --- discharge the same $\AND$ content, encode the witness
--- and the word-level (one packed lane) versus bit-sliced ($\times\wbit$)
asymmetry in how they reach it roughly cancels at the current optimisation
level. The prover is thus a \emph{net} of large offsetting per-phase gaps
(\Cref{app:meas}) rather than a single dominant difference; at the matched
commitment rate it leaves \zinc{} a touch ahead, widening with $N$ as \zinc{}
scales linearly and \binius{} super-linearly.

\subsection{Measured comparison}
\label{sec:results}
\label{sec:res-headline}

All figures are on one Apple~M4 ($10$-core, $16$\,GB), non-zero-knowledge,
provers multithreaded, CPU-only, at the matched commitment rate $1/4$; \zinc{}
is indexed by $\nu=\log_2(\text{rows})$ with ${\approx}(2^\nu-4)/68$
compressions, and \binius{} by block count matched to it. Absolute times are thermally variable
($\sim$$10$--$20\%$ run-to-run), so we read \emph{ratios} and per-phase
\emph{shapes}, which are stable.

\begin{table}[t]
\centering
\renewcommand{\arraystretch}{1.2}
\begin{tabular}{@{}llrrl@{}}
\toprule
& & \binius{} & \zinc{} & winner \\
\midrule
\multirow{3}{*}{$\sim$$1000$ comp}
  & Prove   & $84$\,ms    & $75$\,ms    & \zinc{} $1.1\times$ \\
  & Verify  & $14$\,ms    & $12$\,ms    & \zinc{} $1.2\times$ \\
  & Proof   & $249$\,KiB  & $1.12$\,MiB & \binius{} $4.6\times$ \\
\midrule
\multirow{3}{*}{$\sim$$8000$ comp}
  & Prove   & $1.22$\,s   & $0.58$\,s   & \zinc{} $2.1\times$ \\
  & Verify  & $493$\,ms   & $45$\,ms    & \zinc{} $11\times$ \\
  & Proof   & $299$\,KiB  & $\sim$$2.0$\,MiB & \binius{} $\sim$$7\times$ \\
\bottomrule
\end{tabular}
\caption{Matched rate $1/4$, CPU, at two scales (full sweeps in
\Cref{app:meas}). At moderate $N$ the systems are close on every axis but
proof size; the content is how they \emph{scale}. Representative single runs.}
\label{tab:headline}
\end{table}

\Cref{tab:headline} reads the comparison in two regimes. At $\sim$$1000$
compressions the provers and verifiers are within $\sim$$1.2\times$ and only the
proof differs materially. The structural content is the scaling. \zinc{}'s
prover is linear and \binius{}'s super-linear, so the prover gap opens to
$\sim$$2\times$ by $8000$ compressions and then becomes unbounded: \binius{}'s
$O(N)$ circuit \emph{build} exceeds a minute beyond $\sim$$8000$, where \zinc{}
still proves $15{,}420$ compressions in $\approx 1.2$\,s (and, via swap,
$123{,}361$ on a $16$\,GB host). The verifier and proof move in opposite
directions: \binius{}'s verifier is $O(N)$ --- its in-the-clear wiring sweep
(\Cref{sec:shifts}) --- and \zinc{}'s is $\sqrt N$, dominated by the
$t$-column proximity test, so the verifier ratio grows from near-parity to
$11\times$ and beyond; the proof runs the other way, \binius{} smaller by
$4$--$7\times$ and widening (recursive FRI versus the flat bit-sliced opening).
The per-phase decomposition behind the prover figure, the full sweeps, and the
$O(N)$-versus-$O(1)$ preprocessing data are in \Cref{app:meas}; the mechanism
behind each gap, and when each system wins, in \Cref{sec:analysis,sec:guidance}.

\subsection{The soundness asymmetry}
\label{sec:cmp-soundness}

A fair comparison must record that the two are not at equal soundness status.
\binius{} is fully sound at its $96$-bit target. \zinc{} is unconditionally
sound for the $\AND$ relations and the entire linear part, but its
modular-addition carries are presently \emph{trusted}: the adder operands'
$\psiz$ parents involve a committed carry-out bit and an intra-word $\SHR$ that
do not reduce to row-shift pair-evaluations of committed columns, so $12$ of
the $14$ deployed Hadamards are honest-prover-supplied at the AIR layer
(\Cref{rem:adder-trusted}). The recommended fix replaces those $12$ carry
Hadamards with one grand-product lookup against a public addition table ---
sound by construction, and simultaneously removing the $\times\wbit$ bit-slice
discharge cost for additions. A characteristic-two subtlety is essential: the
lookup must use a \emph{multiplicative} grand product $\prod_i(\gamma-a_i) =
\prod_t(\gamma-t)^{m_t}$, not an additive log-derivative, which is unsound over
an $\Ftwo$-extension because $1/(\gamma-a)+1/(\gamma-a)=0$ lets even
multiplicities cancel. Because closing the gap costs \zinc{} prover time it has
not yet paid, its moderate-$N$ edge includes a small soundness discount; the
asymptotic ($\sim$$2\times$, scaling) win is the more robust claim.

\subsection{When to use which}
\label{sec:guidance-brief}

The three axes are \emph{hash-agnostic} --- they concern the arithmetization,
the proximity test, and the reduction field, not SHA-256 specifically --- so
they yield workload guidance. Hash-dominated work \emph{at scale} favours
\zinc{} on prover throughput, memory, and verifier cost together; a hard
\emph{small standalone proof} favours \binius{} (or a recursion layer wrapped
over \zinc{}, whose $\sqrt N$ verifier is the cheaper inner argument to wrap);
\emph{heterogeneous} statements that pair a hash with big-integer arithmetic
(RSA, ECDSA) favour \zinc{}'s multi-ring composition in a unified limb domain;
and at \emph{moderate} scale the systems are close enough that the binding
output constraint decides. \Cref{sec:guidance} gives the full
workload-by-workload analysis and a decision table.

\section{Conclusion}
\label{sec:conclusion}

\zinc{} and \binius{} agree on the hard decision --- prove SHA-256 in
characteristic two, refusing the prime-field witness blow-up --- and disagree
on three orthogonal downstream choices, each deciding one outcome: uniform AIR
versus word circuit decides the verifier (and the $O(N)$-versus-$O(1)$
preprocessing); flat versus recursive proximity decides the proof; the
base-field versus $\mathrm{GF}(2^{128})$ reduction decides memory and scaling.
The prover, dominated by shared content, is the closest of the four: at the
matched rate $1/4$ on one machine, \zinc{} leads by a touch at moderate $N$,
widening to $\sim$$2\times$ and then unbounded as \binius{}'s $O(N)$ build
walls; its verifier runs from near-parity to past $11\times$ ($\sqrt N$ versus
$O(N)$) on about half the memory, while \binius{}'s recursive proof stays
$4$--$7\times$ smaller, and \binius{} is fully sound where \zinc{}'s adders are
trusted.

Three observations resist the usual summaries. First, \binius{}'s linear
verifier is caused by the cross-entry \emph{wiring}, not the shifts: the
entry-wise shift is succinct in both systems, and the wiring is succinct only
under a uniform arithmetization. Second, the verifier and prover are
\emph{decoupled} in this family --- a succinct verifier costs the prover
nothing to keep and confers nothing on it. Third, the levers each system needs
to close its gap point toward the other: \zinc{}'s proof wants FRI-style
recursion and \binius{}'s prover wants \zinc{}'s free shift, the convergent
endpoint being a binary-field STARK with an $\Ftwo[X]$ cell ring, a uniform
AIR, a base-field discharge, a sound lookup adder, and a recursive proximity
test.
